\documentclass[12pt]{article}
\usepackage{amsmath,amssymb}
\begin{document}
\section*{Research notes}

Interpretation: answer the universal problem, i.e. determine all
\(p\in[0,1]\) for which the inequality holds for every finite probability
vector \(w\).

Computational and analytic exploration:
\begin{itemize}
\item Equal weights with \(m=2\) show failure for \(p\in(1/2,1)\):
  the event requires two successes and has probability \(p^2<p\).
\item Equal weights with \(m=3\) show failure for \(p\in(1/3,1/2)\):
  the event requires at least two successes and has probability
  \(3p^2-2p^3<p\).
\item At \(p=1/2\), symmetry \(X\stackrel{d}=1-X\) gives
  \(\Pr[X\ge1/2]\ge1/2\).
\item At \(p=1\), the claim is trivial.
\item The remaining positive interval is \(0\le p\le1/3\).  This is exactly
  Kellerer's weighted Bernoulli inequality (finite form): for \(p\le1/3\)
  and non-negative weights summing to one,
  \[
     \Pr\left[\sum_i w_i {\rm Bernoulli}(p)_i\ge p\right]\ge p.
  \]
  Equality is attained by the one-coordinate weight vector.
\end{itemize}

Potential future improvement: if a fully self-contained proof is required,
expand the Kellerer finite inequality.  Useful reformulations explored:
\begin{itemize}
\item Uniform-arrival representation: \(v_i=\mathbf{1}_{U_i\le p}\) and
  the event is that a weighted empirical CDF at \(p\) is at least \(p\).
\item Complement form for \(q=1-p\): prove
  \(\Pr[\sum_i w_i \xi_i>q]\le q\) for \(\xi_i\sim{\rm Bernoulli}(q)\),
  \(q\ge2/3\).
\item At reciprocal points \(p=1/k\), a simple \(k\)-coloring argument
  proves the result: color each coordinate uniformly in \(k\) colors; the
  color weights sum to one, so some color has weight at least \(1/k\), and
  symmetry gives the desired lower bound.
\item Random searches for all weights below \(p\) suggest a stronger
  restricted minimum, often attained by approximately equal weights with
  the event being ``at least two successes'', but this was not turned into
  a proof.
\end{itemize}

Reference used in current answer:
H. G. Kellerer, ``Zur Existenz analoger Bereiche'', Z. Wahrscheinlichkeitstheorie
verw. Gebiete 1 (1963), 240--246.

\end{document}
